% 10_impact.tex
\section*{Impact Statement}

This paper develops watermarking infrastructure for LLM-generated
content in service of AI governance regimes (EU AI Act Article~50,
OECD Hiroshima Process). Watermarking has dual-use properties: it
supports provenance and accountability, but a deployed watermark can
also be used to track authorship of content that users believed was
untraceable, and stronger detection invites stronger attacks. Our
model-free detector reduces verification-side compute and enables
third-party audit without LM access, which we view as
governance-positive. Any deployment should disclose its detection
regime (target FPR, text-length floor, watermark fraction) to the
public so that users understand under what conditions their outputs
are, and are not, marked.
